%% ============================================================
%% Comparative Analysis of CNN and Vision Transformer Architectures
%% for Explainable Endometriosis Classification in Laparoscopic Images
%%
%% IEEE format. Compile on Overleaf (recommended) or locally:
%%   pdflatex main  ->  bibtex main  ->  pdflatex main  ->  pdflatex main
%% Switch [journal] to [conference] in \documentclass if targeting a conference.
%% ============================================================
\documentclass[journal]{IEEEtran}

\usepackage{graphicx}
\usepackage{amsmath,amssymb}
\usepackage{booktabs}
\usepackage{multirow}
\usepackage{url}
\usepackage[hidelinks]{hyperref}

\begin{document}

\title{Comparative Analysis of CNN and Vision Transformer Architectures for Explainable Endometriosis Classification in Laparoscopic Images}

\author{
  Author Name\textsuperscript{1},
  Co-Author Name\textsuperscript{2}
  \thanks{\textsuperscript{1}Department, University, City, Country (e-mail: you@example.com).}
  \thanks{\textsuperscript{2}Department, University, City, Country.}
}

\markboth{Preprint --- do not distribute}{}
\maketitle

\begin{abstract}
% TODO: write last (after Results). ~200 words:
% problem, gap, what we did, key numbers, main finding, significance.
\textit{Placeholder abstract. To be written after the Results section.}
\end{abstract}

\begin{IEEEkeywords}
Endometriosis, deep learning, convolutional neural network, vision transformer,
laparoscopy, explainable AI, Grad-CAM, medical image classification.
\end{IEEEkeywords}

% ============================================================
\section{Introduction}
% NOTE: verify every factual number against its cited source before submission.

% ---- Importance of the Research Topic ----
Endometriosis is a chronic, estrogen-dependent inflammatory disorder in which tissue resembling
the endometrial lining grows outside the uterine cavity---on the peritoneum, ovaries, uterus, and,
in advanced cases, as deep infiltrating endometriosis. It affects an estimated 10\% of women of
reproductive age worldwide, translating to roughly 190 million individuals, and is a leading cause
of chronic pelvic pain, dysmenorrhea, dyspareunia, and infertility~\cite{ref_e1,ref_e3}. Beyond
its clinical burden, the condition carries a substantial socioeconomic cost through lost
productivity, repeated consultations, and diminished quality of life, making it a pressing
women's-health research priority that has historically received disproportionately little
attention~\cite{ref_e3,ref_d1}. Improving the speed and accuracy of endometriosis detection
therefore has direct real-world significance for a large, underserved patient population.

% ---- Importance of the Specific Problem ----
Despite its prevalence, endometriosis is exceptionally difficult to diagnose, with patients
experiencing an average diagnostic delay of 7--10 years between symptom onset and definitive
diagnosis~\cite{ref_e3,ref_d1}. This delay arises from non-specific, overlapping symptoms, the
absence of a reliable non-invasive biomarker, and continued reliance on laparoscopy---an invasive
surgical procedure---as the diagnostic gold standard~\cite{ref_a3}. During laparoscopy, accurate
identification of lesions depends heavily on the operating surgeon's experience, producing
inter-observer variability and a real risk of missing subtle or atypical lesions~\cite{ref_a3,ref_a10}.
The consequences of delayed or missed diagnosis are severe: disease progression, avoidable repeat
surgeries, and prolonged suffering~\cite{ref_e4}. An effective, reliable, and objective
decision-support tool to assist lesion recognition during laparoscopy is thus a clear and urgent
clinical need.

% ---- Role of AI / ML / DL in the Domain ----
Artificial intelligence has emerged as a transformative aid for computer-assisted diagnosis in
gynecologic imaging. Convolutional Neural Networks (CNNs) have demonstrated strong performance in
classifying and localizing endometriotic lesions, with reported accuracies of 0.89--0.93 on the
public Gynecologic Laparoscopy ENdometriosis DAtaset (GLENDA)~\cite{ref_a1,ref_a2}, and
meta-analytic evidence confirms deep learning as the strongest-performing family of methods across
imaging modalities~\cite{ref_e2}. More recently, Vision Transformers (ViTs)---which capture global
context through self-attention rather than local convolution---have shown competitive and sometimes
superior performance in surgical image tasks~\cite{ref_a9}. Collectively, these technologies can
reduce observer dependence, standardize interpretation, and shorten the diagnostic pathway,
positioning AI as a meaningful contributor to earlier and more consistent endometriosis diagnosis.

% ---- Existing Limitations and Proposed Contribution ----
Nevertheless, existing image-based studies share several limitations. First, GLENDA is severely
class-imbalanced (373 pathology versus 13{,}438 non-pathology frames), yet many works report
accuracy on imbalanced data, yielding optimistic and potentially misleading results. Second, most
models operate as black boxes, reporting predictive metrics without demonstrating \emph{where} the
model attends or whether that attention aligns with the true lesion, which undermines clinical
trust~\cite{ref_a1,ref_e4}. Third, comparisons between CNN and Transformer families are often
performed without a fair, leakage-free protocol, and frames from the same surgery frequently leak
across training and test sets, inflating reported performance. To address these gaps, this study
proposes a fair, leakage-controlled comparison of CNN and ViT architectures on a balanced GLENDA
subset, augmented with a \emph{quantitative} explainability analysis that measures the overlap
between Grad-CAM attention~\cite{ref_gradcam} and GLENDA's ground-truth lesion masks~\cite{ref_a4}.

\noindent The main contributions of this work are:
\begin{itemize}
  \item A reproducible, leakage-free balanced benchmark on GLENDA for binary endometriosis
        classification, using patient/case-wise splitting so that frames from the same surgery
        never appear in both training and test sets---an issue overlooked in much prior work.
  \item A fair comparison of CNN (ResNet50~\cite{ref_resnet}, EfficientNet-B0~\cite{ref_efficientnet})
        and Vision Transformer (ViT-Base/16~\cite{ref_vit}) architectures under identical
        transfer-learning, augmentation, and cross-validation protocols, evaluated with a complete
        metric suite (accuracy, precision, recall, F1, ROC-AUC) and statistical significance testing.
  \item A quantitative explainability evaluation comparing Grad-CAM attention maps against
        ground-truth lesion masks (IoU-based overlap), assessing not merely \emph{whether} a model
        is accurate but \emph{whether it attends to clinically relevant regions}.
\end{itemize}

These contributions directly address the identified limitations---imbalance, data leakage, and
opaque decision-making---and thereby strengthen the reliability and clinical credibility of
automated endometriosis classification.

The remainder of this paper is organized as follows. Section~\ref{sec:related} reviews related work,
organized by category. Section~\ref{sec:method} describes the dataset and proposed methodology.
Section~\ref{sec:results} presents results and statistical analysis. Section~\ref{sec:discussion}
discusses findings, limitations, and threats to validity, and Section~\ref{sec:conclusion} concludes.

% ============================================================
\section{Literature Review}\label{sec:related}
% TODO (next): overall overview + 2.1 ML-based, 2.2 DL-based, 2.3 XAI, 2.4 Recent works, 2.5 Gap.

\section{Methodology}\label{sec:method}
% TODO

\section{Results}\label{sec:results}
% TODO (after experiments)

\section{Discussion}\label{sec:discussion}
% TODO

\section{Conclusion}\label{sec:conclusion}
% TODO

% ============================================================
\bibliographystyle{IEEEtran}
\bibliography{references}

\end{document}
